\section{Numerical methods and sensitivity analyses}\label{si:numerics}

\paragraph{Numerical method and finite-step implementation check.} The reported two-land simulations use the projected delayed-quality law in Eq.~(10), integrated by the method of steps with explicit first-order Euler updates. At step $k$, current states are read at $t_{k-1}$; the time-dependent technology input and delayed values are evaluated at $t_k$, with linear interpolation of the delayed $q$ and $B$ histories. The raw delayed quality velocity is projected onto $T_{K_q}(q(t_{k-1}))$ before the Euler update. After the state update, $q$ is clipped to $K_q=[0,q_{\max}]$ as a finite-step safeguard; the other bounded states are clipped, reserve area is reconstructed as $A_r=A_{tot}-A_f-A_c$, and the updated $B$ is stored in the delayed history. The positive-part operator is exact. For ordinary runs, $q(s)=q_0$ and $B(s)=B_0$ over the full pre-run interval $[-\max(\tau_g,\tau_p),0]$, with $B_0$ calculated from the initial state. In the technology-wave experiment, both runs share a no-wave prehistory; the active $T_b(0)$ changes current $B(0)$, and later wave-adjusted $B$ values enter $K(t-\tau_p)$. Baseline demonstrations and threshold bisections use $\Delta t=0.05$\,yr; drawdown sweeps use $0.1$\,yr, with selected thresholds rechecked at $0.05$\,yr.

A finite-step consistency test compared post-Euler quality clipping with explicit tangent-cone projection before the Euler update. For the Figure~1 initial state and parameters, $T=100$\,yr, $\Delta t=0.05$\,yr, and $\tau_g=5,20,40$\,yr, the arrays for $A_f$, $A_c$, $q$, $P$, $D$, $A_r$, $B$, $E$, and $T_b$ were exactly equal. For the continuous baseline projected-quality path at $\tau_g=20$ or $40$\,yr, $q(t)=\min(1,0.9+0.0072t)$; its maximum piecewise-linear error fell from $7.1\times10^{-5}$ at $\Delta t=0.1$\,yr to $8.8\times10^{-6}$ at $0.0125$\,yr. In the full frozen-demand Figure~1 scenario, halving the step from $0.05$ to $0.025$\,yr reduced maximum absolute differences against a $0.003125$\,yr reference from $0.00221$ to $0.00102$\,ha for $A_c$, $0.00180$ to $0.000831$\,gha for $B$, $0.000252$ to $0.000118$\,cap for $P$, and $0.00200$ to $0.000924$\,gha\,yr for $D$. These checks establish equality only for the enumerated cases and show smaller differences under timestep refinement; they are not a universal error bound or proof of convergence for all settings.

\paragraph{Selected delayed law and current-quality sensitivity.} The selected model uses $\dot q=\Pi_{T_{K_q}(q(t))}G_c(q(t-\tau_g))$. A current-quality comparator, with the same interval constraint but $G_c(q(t))$, is included only as a sensitivity check and is not the selected model. Under the baseline constant history $q(s)=0.9$, $\rho_c=0.08$\,yr$^{-1}$, and $q_{\max}=1$, the projected delayed-quality paths for $\tau_g=20$ and $40$\,yr are exactly identical: both reach the imposed ceiling at $125/9\approx13.8889$\,yr, before either delay elapses. This is saturation-induced degeneracy, not general delay robustness; the $\tau_g=5$\,yr case is distinct. The current-quality comparator extends the demonstration window by $1.3$\,yr and shifts the floor-proximity threshold by at most $0.001$\,cap in the displayed finite comparison. This finite comparison does not imply equivalence between the two quality laws.

\begin{table}[htbp]
\centering\small
\caption{Current-quality sensitivity comparator versus the selected projected delayed-quality specification. The comparator uses the same imposed interval $K_q$; reported baseline trajectories use the delayed law. The equality of the delayed-quality $\tau_g=20$ and $40$\,yr paths is saturation-induced under the stated history, not evidence of general delay robustness. $B$, $Y_f$, and $\Upsilon(q)A_c$ are in gha; land areas in ha; $P_0$ in cap; delays and step intervals in yr.}
\label{tab:si_delay}
\begin{tabularx}{\textwidth}{>{\raggedright\arraybackslash}Xcc}
\toprule
Quantity & current $q(t)$ & projected delayed $q(t-\tau_g)$ \\
\midrule
Demonstration window with $\dot B>0$, $\dot A_c<0$ (yr) & $150.3$ & $149.0$ \\
Mask-end values $B$, $A_c$, $A_f$ & $0.825,\ 0.808,\ 0.923$ & $0.825,\ 0.808,\ 0.923$ \\
Capital contribution $\Ups(q)A_c$ (initial $\to$ mask end, gha) & $0.08364\to0.04038$ & $0.08364\to0.04038$ \\
Fast-land contribution $Y_f$ (initial $\to$ mask end, gha) & $0.51000\to0.78462$ & $0.51000\to0.78462$ \\
$P_0$ floor threshold ($A_{c0}=1.0$\,ha, $\tau_g=5$\,yr, $\Delta t=0.1$\,yr) & $2.955$ & $2.955$ \\
$P_0$ floor threshold ($A_{c0}=1.0$\,ha, $\tau_g=20$\,yr, $\Delta t=0.1$\,yr) & $2.955$ & $2.954$ \\
$P_0$ floor threshold ($A_{c0}=1.0$\,ha, $\tau_g=40$\,yr, $\Delta t=0.1$\,yr) & $2.955$ & $2.954$ \\
\bottomrule
\end{tabularx}

\end{table}

\paragraph{Well-posedness and scope.} For $\rho_c>0$, $q_{\max}>0$, and an admissible continuous history, the projected quality component is nondecreasing and remains in $K_q$. If $q(0)>0$, it converges to $q_{\max}$; the all-zero history remains zero. The boundary projection makes the vector field constrained and nonsmooth at $q=q_{\max}$, so we do not transfer characteristic roots of the raw unconstrained delay equation to the projected system. The other right-hand sides are locally Lipschitz on each specified active set, and the positive-part and minimum operators are Lipschitz. State bounds and area conservation are enforced as described above. The affine spectrum reported in SI~\S S2 applies only to the one-sided deficit-active land branch; no smooth linear-stability claim is made across the quality ceiling or the shortfall switch.

